% overlap-measure.tex — Chapter on the per-plaquette union bound
% on mu_k(Omega_k^cap) and its summability across RG trajectories.
%
% Linked code / Lean:
%   - computations/overlap_measure.py
%   - lean/YMGap/Sheaf/Site.lean (the configuration site)

\chapter{The Overlap-Measure Union Bound}
\label{ch:overlap-measure}

\noindent\emph{Aggregating the single-plaquette buffer-shell bound
across all plaquettes at scale $k$ via the union bound, and the
structural summability check $\sum_k \muk(\Omkcap) < \infty$.}

This chapter aggregates the single-plaquette bound of
\Cref{ch:buffer-shell} over every plaquette in the scale-$k$ lattice
via the standard union bound, producing a closed-form upper bound on
the buffer-region measure $\muk(\Omkcap)$.  The aggregate bound has
the same $(C, \kappa, \beta)$ shape as its per-plaquette parent,
multiplied by the plaquette count $N_p(k)$; summability of the
sequence $\{\muk(\Omkcap)\}_k$ across an RG trajectory is the
structural condition that makes the buffer-shell sequence harmless
for the inductive-step closure.  The chapter is short and
classical; its content is structural, not analytical.

\section{The buffer event at scale \texorpdfstring{$k$}{k}}
\label{sec:overlap-setup}

Let $E_k$ be the set of (oriented) lattice edges at scale $k$ and
$\Omk = \SUtwo^{E_k}$ the link-configuration space.  Let $\Pcal_k$
be the set of plaquettes (oriented two-dimensional faces of the
hypercubic lattice at scale $k$), with cardinality $N_p(k)$.  For
each plaquette $p \in \Pcal_k$, write $U_p \in \SUtwo$ for the
plaquette holonomy (the oriented product of link variables around
$p$).

\marginkey{Buffer region $\Omkcap$}%

\begin{definition}[Per-plaquette buffer event and buffer region]
\label{def:buffer-region}
Fix the small-field smallness parameter $\delta_k > 0$ as in
\Cref{ch:buffer-shell}, Definition~\ref{def:buffer-annulus}.  The
per-plaquette buffer event is
\begin{equation}
  \label{eq:per-plaquette-event}
  A_p \;:=\; \bigl\{ U \in \Omk : U_p \in A(\delta_k) \bigr\},
\end{equation}
and the (aggregate) \emph{buffer region} at scale $k$ is
\begin{equation}
  \label{eq:buffer-region}
  \Omkcap \;:=\; \bigcup_{p \in \Pcal_k} A_p \;\subseteq\; \Omk .
\end{equation}
$\Omkcap$ is the subset of configurations where at least one
plaquette holonomy lands in the buffer annulus.
\end{definition}

The buffer region is the dynamical setting for the sheaf-theoretic
overlap analysis of \Cref{ch:sheaf-framework} (the three-element
cover $\{\Omksf, \Omklf, \Omkcap\}$).  We need an upper bound on
its Gibbs measure $\muk(\Omkcap)$.

\section{The union bound}
\label{sec:union-bound}

\marginkey{The union bound}%

\begin{theorem}[Per-plaquette union bound]
\label{thm:union-bound}
Under the heat-kernel Gibbs measure $\muk$,
\begin{equation}
  \label{eq:union-bound}
  \muk(\Omkcap)
  \;=\;
  \muk\!\left(\bigcup_{p \in \Pcal_k} A_p\right)
  \;\le\;
  \sum_{p \in \Pcal_k} \muk(A_p)
  \;=\;
  N_p(k) \cdot P_{\mathrm{shell}}(\delta_k, g_k),
\end{equation}
where the last equality uses translation invariance of $\muk$
(every $A_p$ has the same probability $P_{\mathrm{shell}}(\delta_k,
g_k)$ of \Cref{ch:buffer-shell},
equation~\eqref{eq:buffer-shell-prob}).
\end{theorem}

\begin{proof}
The first inequality is the union bound for the probability of a
finite union of events.  The equality uses that the heat-kernel
Gibbs measure on $\Omk$ is invariant under lattice translations, so
$\muk(A_p)$ is independent of $p$.
\end{proof}

Combining \cref{thm:union-bound} with the single-plaquette
uniform-in-$g$ bound of \Cref{ch:buffer-shell},
Theorem~\ref{thm:buffer-uniform}:

\begin{corollary}[Aggregate buffer-region bound]
\label{cor:aggregate-bound}
For $\delta_k \in (0.02, 1.0)$ and $g_k \in (0.05, 1.5)$,
\begin{equation}
  \label{eq:aggregate-bound}
  \muk(\Omkcap)
  \;\le\;
  N_p(k) \cdot \bigl(\, P_{\mathrm{Haar}}(\delta_k)
    \;+\; C_{\mathrm{tail}} \cdot \delta_k^3 \cdot \exp(-\kappa(\SUtwo) / g_k^2)\,\bigr) ,
\end{equation}
with $\kappa(\SUtwo) = \tfrac34$ in the canonical-$\kappa$
convention and $C_{\mathrm{tail}} = 4$.
\end{corollary}

\section{Summability across an RG trajectory}
\label{sec:overlap-summability}

The structural question raised by the buffer construction is:
\emph{does the buffer-region measure remain controlled across the
RG trajectory?}  Concretely, given a trajectory of couplings
$\{g_k\}_{k \ge 0}$ and smallness parameters $\{\delta_k\}_{k \ge 0}$,
is the sequence $\{\muk(\Omkcap)\}_k$ summable?

\marginkey{Geometric plaquette count}%

\begin{lemma}[Plaquette count under block-spinning]
\label{lem:plaquette-count}
Under block-spin coarsening with block factor $b > 1$ in four
dimensions, $a_{k+1} = b\, a_k$, so $N_p(k+1) = b^{-4}\, N_p(k)$.
Hence $N_p(k) = N_p(0)\, b^{-4k}$.
\end{lemma}

\begin{proof}
Direct: the number of plaquettes in a volume of fixed physical size
$L^4$ at lattice spacing $a_k$ is approximately $6\,(L/a_k)^4$;
block-spinning replaces $a_k$ by $b\,a_k$, dividing the count by
$b^4$.
\end{proof}

\begin{theorem}[Per-step bound + summability sufficient condition]
\label{thm:union-summability}
Let $\{(g_k, \delta_k)\}_{k \ge 0}$ be an RG trajectory satisfying:
\begin{enumerate}[label=(\arabic*)]
  \item the buffer-shell bound of \Cref{ch:buffer-shell},
    Theorem~\ref{thm:buffer-uniform}, applies at every $k$ (i.e.,
    $(\delta_k, g_k) \in (0.02, 1.0) \times (0.05, 1.5)$);
  \item the smallness parameter satisfies $\delta_k \to 0$ as
    $k \to \infty$ (small-field convergence in the UV-direction
    trajectory).
\end{enumerate}
Then
\begin{equation}
  \label{eq:union-summability}
  \sum_{k \ge 0} \muk(\Omkcap)
  \;\le\;
  N_p(0) \sum_{k \ge 0} b^{-4k}\,
    \bigl(\,P_{\mathrm{Haar}}(\delta_k) + 4\,\delta_k^3\,e^{-3/(4 g_k^2)}\,\bigr)
  \;<\; \infty,
\end{equation}
the right-hand side being finite by \cref{lem:plaquette-count} and
the geometric factor $b^{-4k}$.
\end{theorem}

\begin{proof}[Proof sketch]
The per-step bound is \cref{cor:aggregate-bound} multiplied by
$N_p(k) = N_p(0)\,b^{-4k}$ from \cref{lem:plaquette-count}.  The
sum is bounded above by a geometric series with ratio $b^{-4}$,
which is convergent for any block factor $b > 1$, and the per-step
factor $(P_{\mathrm{Haar}}(\delta_k) + 4\,\delta_k^3\,e^{-\kappa/g_k^2})$
is bounded uniformly in $k$ on the convergent trajectory.
\end{proof}

\marginkey{Structural, not analytical}%

\Cref{thm:union-summability} is a \emph{structural} closure: it
guarantees the aggregate buffer-region measure remains controlled
across the trajectory.  The stronger \emph{analytical} closure —
that the $\sigh$-budget $\sum_k \eta_k$ stays below $\sigma_0$ — is
the content of \Cref{ch:sigma-hat-trajectory}; the buffer-region
summability of this chapter is a necessary precursor to that
closure, not a substitute.

\section{Looseness of the union bound}
\label{sec:union-looseness}

The per-plaquette events $\{A_p\}_{p \in \Pcal_k}$ are
\emph{positively correlated}: adjacent plaquettes share an edge, so
the Wilson-action heat-kernel measure conditions $U_p \in A(\delta_k)$
into a probability of $U_{p'} \in A(\delta_k)$ that is greater than
the unconditional value, for any face-neighbor plaquette $p'$.  The
union bound's worst-case independence assumption is exactly wrong
in this direction: the true measure $\muk(\Omkcap)$ is strictly
less than the right-hand side of \eqref{eq:union-bound}, and the
overestimation factor grows with the correlation length $\xi(g)$
between plaquettes.

\begin{definition}[Union-bound looseness factor]
\label{def:looseness}
At scale $k$ with parameters $(\delta_k, g_k)$, the union-bound
looseness factor is
\begin{equation}
  \label{eq:looseness-factor}
  L(\delta_k, g_k) \;:=\; \frac{N_p(k) \cdot P_{\mathrm{shell}}(\delta_k, g_k)}{\muk(\Omkcap)}
  \;\in\; [1, \infty),
\end{equation}
with $L = 1$ in the strictly-independent (Gaussian) limit and
$L > 1$ at any nontrivial correlation length.
\end{definition}

\marginkey{Conjecture: $L \le 2$}%

\begin{conjecture}[Bounded looseness]
\label{conj:gaussian-sharpness}
There exists a universal constant $C \le 2$ such that for every
trajectory and every $k$ in the convergent regime,
$L(\delta_k, g_k) \le C$.  Empirically, across the canonical $\SUtwo$
trajectory we observe $L \le 1.18$.
\end{conjecture}

The looseness conjecture motivates \Cref{ch:chaining-estimate},
where a correlation-aware tightening replaces the union bound with
a chaining estimate that captures the positive-correlation
reduction explicitly.  In the regime where the chaining estimate
is admissible (small-positive correlations, $g_k \le 2$), the
chaining bound is strictly tighter than the union bound; outside
that regime the union bound remains in use.

\section{Code provenance}
\label{sec:overlap-codeanchor}

\begin{itemize}
  \item Python implementation:
    \codepath{computations/overlap\_measure.py}.  Key entry points:
    \codepath{union\_bound\_overlap(n\_plaquettes, g\_k, delta\_k,
    t\_choice)} (\cref{thm:union-bound}, returning
    \codepath{UnionBoundResult}); \codepath{is\_summable\_in\_k}
    (\cref{thm:union-summability}); and the standard SU(2)
    trajectory generator
    \codepath{standard\_su2\_trajectory(k\_max, g\_0)}.  Tests:
    \codepath{tests/test\_overlap\_measure.py}.
  \item Visualizations:
    \codepath{visualizations/pilot/su2/overlap-vs-rg-step.svg}
    plots $\muk(\Omkcap)$ across the canonical trajectory with the
    union bound;
    \codepath{visualizations/pilot/su2/conjecture-7.4-looseness.svg}
    quantifies $L(\delta_k, g_k)$ empirically.
  \item The chaining tightening of \cref{conj:gaussian-sharpness}
    is in \codepath{computations/chaining\_estimate.py} and is the
    subject of \Cref{ch:chaining-estimate}.
\end{itemize}
